\chapter{Исследовательская часть}

\section{Оценка временной эффективности}
\subsection{Поиск полным перебором}
Лучший случай --- искомый элемент является первым элементом массива. В данном случае будет выполнено одно сравнение искомого элемента с первым элементом массива.

Худший случай --- искомого значения в массиве нет или искомый элемент является последним в массиве. В данной ситуации необходимо будет проверить все элементы массива, то есть выполнить $N$ сравнений.

\subsection{Бинарный поиск}
В реализации бинарного поиска искомый элемент сравнивается со средним элементом, после чего размер массива сокращается.

Пусть $N_i$~--- размер рассматриваемой части массива после $i$-й итерации, где $N_0 = N$~--- исходный размер массива. Тогда
\begin{align}
	N_1 &= \frac{N_0}{2}, \notag \\
	N_2 &= \frac{N_1}{2} = \frac{N_0}{4}, \notag \\
	&\;\;\vdots \notag \\
	N_i &= \frac{N_0}{2^i}.
\end{align}
Поиск завершается, когда в рассматриваемой части остаётся один элемент: $N_0 / 2^i = 1$, откуда $i = \log_2 N$.
Уменьшение массива в два раза на каждой итерации позволяет уменьшить количество итераций до $\log_2N$. Число сравнений зависит от количества итераций и возможности предварительного выхода.

\subsubsection{Нерекурсивная реализация без предварительного выхода}
Для нерекурсивной реализации без предварительного выхода лучший и худший случаи будут совпадать, поскольку границы будут сужаться, пока не будут равны. В итоге для определения индекса необходимо $\log_2N + 1$ сравнений --- поиск индекса в массиве и сравнение искомого значения с элементом по найденному индексу.

\subsubsection{Нерекурсивная реализация с предварительным выходом}
В нерекурсивной реализации с предварительным выходом на каждой итерации выполняется сравнение текущего среднего элемента с искомым. Если в результате сравнения найден искомый элемент, то процесс поиска завершается.

Лучший случай --- искомый элемент находится в середине массива. В первой же итерации найдём индекс искомого элемента в массиве и завершим поиск. В итоге необходимо только одно сравнение среднего элемента с искомым.

Худший случай --- искомый элемент на границах массива или его нет в массиве. На каждой итерации будет выполняться проверка равенства среднего элемента с искомым и проверка того, что средний элемент меньше искомого. Таким образом, получим $2\log_2N$ сравнений.

\subsubsection{Рекурсивная реализация}
Рекурсивная реализация по количеству сравнений полностью совпадает с нерекурсивной реализацией с предварительным выходом.

\section{Оценка затрачиваемой памяти}
При оценке затрачиваемой памяти учитываются только затраты сверх входных данных.

\subsection{Поиск полным перебором}
Поиск полным перебором использует две переменные:
\begin{itemize}
	\item $i$ --- индекс текущего рассматриваемого элемента;
	\item $index$ --- индекс искомого элемента.
\end{itemize}
Вызов функции производится один раз. 
Дополнительные затраты составляют константное значение --- $O(1)$.

\subsection{Нерекурсивный бинарный поиск}
Нерекурсивный бинарный поиск использует три переменные:
\begin{itemize}
	\item $l$ --- индекс левой границы;
	\item $r$ --- индекс правой границы;
	\item $m$ --- индекс среднего элемента.
\end{itemize}
Вызов функции производится один раз.
Дополнительные затраты составляют константное значение --- $O(1)$.

\subsection{Рекурсивный бинарный поиск}
Пусть $M_s$ --- память, затрачиваемая на один кадр стека; он включает в себя: $array$ (ссылка), $t$, $l$, $r$, $m$. 
Рекурсивный бинарный поиск использует одну переменную:
\begin{itemize}
	\item $m$ --- индекс среднего элемента.
\end{itemize}
Глубина рекурсии составляет $\log_2N$.
Дополнительные затраты составляют $\log_2N \cdot M_s$ --- $O(\log_2N)$.

\section{Сравнительный анализ}
Результаты сравнительного анализа приведены в таблице~\ref{tbl:search-comparison}, где ЛС~--- лучший случай, ХС~--- худший случай.
\input{parts/comparison-table.tex}
\section{Вывод}
Алгоритм поиска полным перебором в худшем случае выполняет $N$ сравнений искомого элемента с элементами массива, при этом не налагает дополнительных требований на массив. 

Алгоритмы бинарного поиска позволяют сократить количество сравнений до логарифмического, но требуют поддержания массива в отсортированном состоянии. Использование предварительного выхода в бинарном поиске позволяет в лучшем случае сократить количество сравнений до одного, но в худшем случае число сравнений возрастает в два раза. Временная эффективность рекурсивной реализации совпадает с эффективностью нерекурсивной реализации с предварительным выходом, но рекурсивная реализация дополнительно затрачивает память на кадры стека.

Для однократного поиска в неотсортированном массиве рекомендуется поиск полным перебором, для отсортированного массива~--- нерекурсивный бинарный поиск без предварительного выхода как наиболее экономный по сравнениям в худшем случае и по памяти.